\chapter*{Заключение}
\addcontentsline{toc}{chapter}{Заключение}

В ходе работы рассмотрено применение методов машинного обучения и больших языковых моделей для анализа малого производственного бизнеса. На примере компании по производству кондитерских изделий показано, что банковские и операционные данные могут быть преобразованы в инструмент управленческого контроля.

Разработанный подход включает очистку исходных выгрузок, классификацию банковских операций, сравнение rule-based и LLM-разметки, построение управленческого cash flow, финансовый аудит и обучение прогнозных моделей.

Главный вывод финансового анализа состоит в том, что компания имеет существенный оборот, но почти не формирует операционный денежный поток. Операционные поступления и списания практически совпадают, а денежная операционная маржа составляет около 0,05\%. Это объясняет, почему бизнес может работать и выполнять обязательства, но не давать стабильный свободный доход собственникам.

Практическая ценность ML и LLM в данной задаче проявляется на трех уровнях. Во-первых, автоматическая разметка превращает банковскую выписку в управленческую таблицу. Во-вторых, финансовый аудит выявляет зоны риска: производственные закупки, кассовые провалы, концентрацию поставщиков и зависимость от ключевого канала продаж. В-третьих, прогнозные модели создают основу для регулярного управления спросом, возвратами и cash flow.

Ограничения работы связаны с короткой историей данных, неполными справочниками, необходимостью ручной валидации и чувствительностью моделей к крупным разовым платежам. Дальнейшее развитие проекта предполагает добавление фактических остатков на счетах, календаря обязательных платежей, расширенного справочника себестоимости, истории запусков моделей и production-loop переобучения.
